\section{Conclusion}
\label{sec:conclusion}

K-means clustering of TIGER fingerprint space identifies four sparse clusters in
the 20-game validation corpus, three of which are validated by a five-designer expert
panel as genuinely unexplored and commercially viable design opportunities, surpassing
the pre-registered threshold of $\geq 2$ validated gaps. This demonstrates that TIGER is not only a framework for evaluating
existing games but a generative tool for discovering where new games could be made.

The broader contribution is a methodology for computational design space exploration:
score a corpus, normalize to remove quality bias, project into a low-dimensional
framework space, identify sparse regions, and translate sparse coordinates into
natural-language design briefs for expert validation. This pipeline is domain-agnostic
and could be applied to other creative artifact corpora (video games, film narratives,
musical compositions) where multi-attribute scoring data is available.

TIGER profiles for all 150 corpus games, clustering code, and expert validation
transcripts are available at [URL]. A full appendix listing TIGER coordinates for
all 150 games is provided in the supplementary materials.
